\textbf{\textit{Soal. }}An acute-angled triangle $ABC$ is fixed on a plane with largest side $BC$. Let $PQ$ be an arbitrary diameter of its circumscribed circle, and the point $P$ lies on the smaller arc $AB$, and the point $Q$ is on the smaller arc $AC$. Points $X, Y, Z$ are feet of perpendiculars dropped from point $P$ to the line $AB$, from point $Q$ to the line $AC$ and from point $A$ to line $PQ$. Prove that the center of the circumscribed circle of triangle $XYZ$ lies on a fixed circle.

\textit{Terjemahan: } Diberikan sebuah segitiga lancip $ABC$ yang tetap di sebuah bidang dengan sisi terpanjangnya adalah $BC$. Misalkan $PQ$ adalah sebuah diameter sembarang dari lingkaran luar segitiga $ABC$, dengan titik $P$ berada di busur $AB$ yang lebih kecil  dan titik $Q$ berada di busur $AC$ yang lebih kecil. Titik $X,Y,Z$ berturut-turut adalah kaki tinggi titik $P$ ke garis $AB$, titik $Q$ ke garis $AC$, dan titik $A$ ke garis $PQ$. Buktikan bahwa pusat lingkaran luar segitiga $XYZ$ berada pada sebuah lingkaran tetap.

\begin{proof}[\textbf{Solusi. }] (Thursday, 27 October 2022) Misalkan $l_1$ dan $l_2$ berturut-turut adalah garis sumbu $XZ$ dan $ZY$. Misalkan $l_1 \cap AB = M$, $l_1 \cap XZ = K$, $l_2 \cap ZY = L$ dan $l_2 \cap AC = N$.  Terakhir definisikan $O_1, O_P, O_Q$ berturut-turut sebagai pusat lingkaran $(XYZ)$, $(PXA)$, $(QYA)$. Jelas $O_1$ merupakan perpotongan $l_1$ dan $l_2$, $l_1$ melewati $O_P$ dan $l_2$ melewati $O_Q$.

% \begin{center}
    \begin{asy}
        import olympiad;
        import geometry;
        size(350);
        
        pair A, B, C;
        A = dir(125);
        B = dir(-150);
        C = dir(-30);
        dot("$A$", A, NW);
        dot("$B$", B, SW);
        dot("$C$", C, SE);
        draw(A--B--C--A);
        circle c = circumcircle(triangle(A,B,C));
        draw(c);
        
        pair O, P, Q;
        O = circumcenter(triangle(A,B,C));
        P = dir(187);
        Q = rotate(180, O)*P;
        dot("$O$", O, NW);
        dot("$P$", P, W);
        dot("$Q$", Q, E);
        draw(P--Q);

        pair X, Y, Z;
        X = foot(P, A, B);
        Y = foot(Q, A, C);
        Z = foot(A, P, Q);
        dot("$X$", X, SE);
        dot("$Y$", Y, N);
        dot("$Z$", Z, S);
        draw(P--X, blue);
        draw(Q--Y, blue);
        draw(A--Z, blue);
        draw(X--Y--Z--X, red);
        draw(circumcircle(triangle(X,Y,Z)), blue);
        pair O1 = circumcenter(X,Y,Z);
        dot("$O_1$", O1, S);
        
        pair K,L,M,N, Op, Oq;
        K = (X+Z)/2;
        L = (Z+Y)/2;
        M = intersectionpoint(line(O1, K), line(A,B));
        Op = intersectionpoint(line(O1, K), line(A,P));
        N = intersectionpoint(line(O1, L), line(A,C));
        Oq = intersectionpoint(line(O1, L), line(A,Q));
        dot("$K$", K, NW);
        dot("$L$", L, NE);
        dot("$M$", M, NW);
        dot("$N$", N, NE);
        dot("$O_P$", Op, NW);
        dot("$O_Q$", Oq, NE);
        
        draw(O1--Op, green);
        draw(O1--Oq, green);
        draw(O1--X, green);
        draw(O1--Y, green);
        draw(O1--Z, green);
        draw(A--P, blue);
        draw(A--Q, blue);
        draw(P--B, green);
        draw(Q--C, green);
        draw(circumcircle(triangle(K,L,O1)), purple);
        draw(circumcircle(triangle(M,N,O1)), dashed+red);
        draw(circumcircle(triangle(P,X,Z)), purple);
        draw(circumcircle(triangle(Q,Y,Z)), purple);

        draw(rightanglemark(A, Z, Q, 2));
        draw(rightanglemark(P, X, B, 2));
        draw(rightanglemark(Q, Y, C, 2));
        draw(rightanglemark(Op, K, X, 2));
        draw(rightanglemark(Oq, L, Y, 2));
    \end{asy}
\end{center}
\begin{center}
    \begin{tikzpicture}[scale=4, dot/.style={circle, fill, inner sep=0.5pt, outer sep=0pt}]

        \tkzDefPoint(125:1){A}
        \tkzDefPoint(-150:1){B}
        \tkzDefPoint(-30:1){C}
        \tkzDefTriangleCenter[circum](A,B,C)\tkzGetPoint{O}

        % P pada lingkaran luar, Q titik antipodalnya
        \tkzDefPoint(187:1){P}
        \tkzDefPointBy[symmetry=center O](P)\tkzGetPoint{Q}

        % Kaki-kaki tegak lurus
        \tkzDefPointBy[projection=onto A--B](P)\tkzGetPoint{X}
        \tkzDefPointBy[projection=onto A--C](Q)\tkzGetPoint{Y}
        \tkzDefPointBy[projection=onto P--Q](A)\tkzGetPoint{Z}

        \tkzDefTriangleCenter[circum](X,Y,Z)\tkzGetPoint{O1}

        \tkzDefMidPoint(X,Z)\tkzGetPoint{K}
        \tkzDefMidPoint(Z,Y)\tkzGetPoint{L}
        \tkzInterLL(O1,K)(A,B)\tkzGetPoint{M}
        \tkzInterLL(O1,K)(A,P)\tkzGetPoint{Op}
        \tkzInterLL(O1,L)(A,C)\tkzGetPoint{N}
        \tkzInterLL(O1,L)(A,Q)\tkzGetPoint{Oq}

        \tkzDrawPolygon(A,B,C)
        \tkzDrawCircle(O,A)
        \tkzDrawSegment(P,Q)
        \tkzDrawSegment[blue](P,X)
        \tkzDrawSegment[blue](Q,Y)
        \tkzDrawSegment[blue](A,Z)
        \tkzDrawPolygon[red](X,Y,Z)
        \tkzDrawCircle[blue](O1,X)
        \tkzDrawSegments[green](O1,Op O1,Oq O1,X O1,Y O1,Z)
        \tkzDrawSegments[blue](A,P A,Q)
        \tkzDrawSegments[green](P,B Q,C)
        \tkzDefTriangleCenter[circum](K,L,O1)\tkzGetPoint{Oklo}\tkzDrawCircle[purple](Oklo,K)
        \tkzDefTriangleCenter[circum](M,N,O1)\tkzGetPoint{Omno}\tkzDrawCircle[dashed,red](Omno,M)
        \tkzDefTriangleCenter[circum](P,X,Z)\tkzGetPoint{Opxz}\tkzDrawCircle[purple](Opxz,P)
        \tkzDefTriangleCenter[circum](Q,Y,Z)\tkzGetPoint{Oqyz}\tkzDrawCircle[purple](Oqyz,Q)

        \tkzMarkRightAngle[size=0.1](A,Z,Q)
        \tkzMarkRightAngle[size=0.1](P,X,B)
        \tkzMarkRightAngle[size=0.1](Q,Y,C)
        \tkzMarkRightAngle[size=0.1](Op,K,X)
        \tkzMarkRightAngle[size=0.1](Oq,L,Y)

        \tkzDrawPoints(A,B,C,O,P,Q,X,Y,Z,O1,K,L,M,N,Op,Oq)
        \tkzLabelPoint[above left](A){$A$}
        \tkzLabelPoint[below left](B){$B$}
        \tkzLabelPoint[below right](C){$C$}
        \tkzLabelPoint[above left](O){$O$}
        \tkzLabelPoint[left](P){$P$}
        \tkzLabelPoint[right](Q){$Q$}
        \tkzLabelPoint[below right](X){$X$}
        \tkzLabelPoint[above](Y){$Y$}
        \tkzLabelPoint[below](Z){$Z$}
        \tkzLabelPoint[below](O1){$O_1$}
        \tkzLabelPoint[above left](K){$K$}
        \tkzLabelPoint[above right](L){$L$}
        \tkzLabelPoint[above left](M){$M$}
        \tkzLabelPoint[above right](N){$N$}
        \tkzLabelPoint[above left](Op){$O_P$}
        \tkzLabelPoint[above right](Oq){$O_Q$}

    \end{tikzpicture}
\end{center}


Perhatikan bahwa $\dangle PXA = \dangle PZA$ menyebabkan $P,X,Z,A$ konsiklis. Dengan cara yang sama $Q,Y,Z,A$ konsiklis. Di lain sisi, ingat bahwa $PQ$ adalah diameter lingkaran $(APBCQ)$ serta $NL \perp ZY$. Sehingga kita punya $$\dangle O_QNA = \dangle LNY = 90^\circ-\dangle NYL = 90^\circ-\dangle AYZ = 90^\circ-\dangle AQZ = 90^\circ-\dangle AQP=\dangle QPA = \dangle QCA$$ yang menyebabkan $O_QN \parallel QC$. Dengan cara yang sama diperoleh juga $O_PM \parallel PB$. Karena $O_Q$ dan $O_P$ masing-amsing merupakan titik tengah $AQ$ dan $AP$, maka $M$ dan $N$ masing-masing merupakan titik tengah $AB$ dan $AC$. Hal ini menyebabkan $l_1$ dan $l_2$ berturut-turut pasti melewati titik tengah $AB$ dan $AC$ untuk seluruh kemungkinan letak $P,Q$. Oleh karena itu, cukup dibuktikan bahwa besar $\dangle MO_1N$ yang akan membuktikan bahwa lingkaran yang melewati $M,N,O_1$ tetap.

Perhatikan karena $\dangle ZXO_1 = \dangle ZLO_1 = 90^\circ$ maka $XZLO_1$ siklis. Oleh karena itu 
\begin{align*}
    \dangle MO_1N 
    &= \dangle KO_1L\\ 
    &= \dangle KZL\\
    &= \dangle XZA + \dangle AZY\\ 
    &= (\dangle XAZ + \dangle ZXA) + (\dangle ZAY + \dangle AYZ)\\
    &= (\dangle XAZ + \dangle ZAY)+ (\dangle ZXA + \dangle AYZ)\\
    &= \dangle BAC + (\dangle ZPA + \dangle AQZ)\\
    &= \dangle BAC + 90^\circ
\end{align*}

Karena $\dangle BAC$ tetap, maka $\dangle MO_1N$ tetap, sesuai yang ingin dibuktikan.
    
\end{proof}